% appendix_reproducibility.tex — Appendix B
% Describes the companion GitHub repository and the reproducibility /
% provenance / machine-readability ideas it explores. Cited from the
% introduction; introduces no new scientific claims.

\section{The companion repository: an experiment in machine-readable provenance}
\label{app:reproducibility}

The introduction notes that all code, data-loading scripts, and the analysis
pipeline are public.\footnote{\url{https://github.com/matiaszaldarriaga/desi-w0wa-svd}}
This appendix describes what the repository contains and, more importantly, the
ideas it is meant to explore. The science in this paper is modest; the part I
think is worth a reader's attention is the \emph{form} in which it is released.
The repository is an experiment in using an AI coding agent not only to do the
calculation but to package the resulting paper so that every number, figure, and
claim is traceable to the code that produced it, and so that the whole argument
is legible to \emph{other} agents as well as to people. I make no claim that the
specific conventions below are the right ones; they are offered as a concrete
example that readers --- and their agents --- can inspect, criticize, and improve.

\subsection{What is in the repository}
\label{app:repo_contents}

Beyond the usual contents (the LaTeX source, the analysis package, and the
scripts that turn the public DESI~DR2 BAO, supernova, and CMB MCMC chains into
the distance-ratio decompositions used here), the repository is organized around
three goals: exact reproducibility, end-to-end provenance of every quoted number,
and a machine-readable map of the paper's claims and evidence. A top-level
\texttt{README} documents the layout and the commands below.

\subsection{Reproducibility}
\label{app:reproducibility_pipeline}

The pipeline is deliberately split so that the expensive, data-heavy step is
separated from figure generation. One script loops over the MCMC chain samples
once, computing the cosmological distances, whitened covariances, SVD
decompositions, and grid evaluations, and caches every resulting array. A second
script computes every number quoted in the paper from those arrays. A third --- a
marimo notebook that also runs as plain Python --- reads only the cached arrays
and regenerates every figure in seconds. Because the cached arrays and the number
file are committed, a reader who only wants to rebuild the paper need not download
the multi-gigabyte chains at all; the full pipeline from the public chains is a
single orchestration script for anyone who wants to verify the cache itself. The
intent is that the gap between ``download the inputs'' and ``reproduce every
figure'' is a documented command, not an exercise left to the reader.

\subsection{Provenance of every number}
\label{app:number_provenance}

Every number in the paper text carries a \texttt{\textbackslash dataref}
annotation in the LaTeX source that names a key. That key resolves to an entry in
a committed machine-readable file (\texttt{paper\_numbers.json}), and the same key
appears in the script that computes it. A reader who wants to know where, say, the
$+2.2\sigma$ BAO \czero tension comes from can follow the key from the printed
sentence to the stored value to the exact line of code and the formula it
evaluates. Nothing in the paper is a hand-typed number: the build regenerates the
value file, so a stale or mistyped figure in the text is detectable rather than
silent. This is the provenance idea I would most like readers to evaluate ---
whether tying each printed number to a key, a stored value, and a line of code is
worth the modest bookkeeping it costs.

\subsection{Claims, evidence, and the dependency graph}
\label{app:claims_structure}

The repository carries a set of machine-readable registries, maintained alongside
the LaTeX, that record the logical skeleton of the paper rather than its prose. A
claim registry lists each scientific claim with a stable identifier, its statement,
pointers to the figures, stored numbers, and code that constitute its evidence, and
links to the other claims it depends on. Companion registries do the same for
figures (each tied to the script and notebook cell that produces it and the claims
it supports) and for the critical-path scripts. A separate argument-skeleton file
records what each section establishes and what it rests on, and a dependency graph
renders the claim-to-claim structure directly. The aim is that the paper's argument
exists not only as text to be read linearly but as an explicit graph that can be
queried: \emph{which claims depend on this result?}, \emph{what is the evidence for
that one?}, \emph{which figure and which line of code support it?}

\subsection{An interactive view, and an agent-facing one}
\label{app:dashboard}

The same registries drive an interactive HTML dashboard that renders the
claim--evidence graph in a browser, so a reader can navigate from any claim down to
the figures, numbers, and code that support it without reading the source files.
The registries are equally meant to be consumed by software: the annotation macros
and the YAML files were built so that an agent handed this repository can locate
every claim, verify that each traces to code and data, and check the paper for
internal consistency --- the same verification I relied on while writing it. The
broader concept behind the repository is that as papers are increasingly read,
checked, and built upon by AI agents, it is worth publishing the structure of an
argument in a form those agents can ingest directly, rather than leaving them to
reconstruct it from prose. The repository is offered as a worked example of that
idea.
